%!TEX root = ./main.tex

\section[Foundation]{The 802.11 Foundation}

% SECTION TIMING: 9:00--17:00 (2 frames). No slack: if running behind, the donor is
% the "Original PHY choices" column on the 1997 frame -- name the three PHYs and move on.

% Teaching note (240 s): this is the deck's load-bearing diagram. The MAC students
% recognize -- association, ACKs, CSMA/CA -- has survived since 1997; everything below
% it has been rewritten repeatedly. Define the PPDU here, once, because every later
% slide quietly assumes it: the PHY-layer unit actually placed on air. When a later
% slide says "OFDMA packs four users into one PPDU", this is the box it means.
\begin{frame}{The 802.11 stack: stable MAC, moving PHY}
  \centering
  % SOURCE: https://standards.ieee.org/ieee/802.11/10548/
  \begin{tikzpicture}[scale=0.80,transform shape]
    \node[draw=softgray,fill=gray!8,rounded corners=3pt,minimum width=8.2cm,
          minimum height=0.85cm,align=center] (llc) at (0,0)
      {\textbf{LLC / IP / applications}\\[-2pt]{\scriptsize outside 802.11}};
    \node[netnode,minimum width=8.2cm,minimum height=0.95cm,below=0.25cm of llc] (mac)
      {802.11 MAC\\[-2pt]{\scriptsize\mdseries association \,\textperiodcentered\,
        framing \,\textperiodcentered\, ACK \,\textperiodcentered\, security
        \,\textperiodcentered\, CSMA/CA}};
    \node[netnode,minimum width=8.2cm,minimum height=0.95cm,below=0.25cm of mac] (plcp)
      {PLCP\\[-2pt]{\scriptsize\mdseries preamble \,\textperiodcentered\, header
        \,\textperiodcentered\, rate signaling \,\textperiodcentered\, PPDU format}};
    \node[netnode,minimum width=8.2cm,minimum height=0.95cm,below=0.25cm of plcp] (pmd)
      {PMD / waveform\\[-2pt]{\scriptsize\mdseries DSSS \,\textperiodcentered\, OFDM
        \,\textperiodcentered\, MIMO \,\textperiodcentered\, OFDMA
        \,\textperiodcentered\, single carrier}};
    \draw[decorate,decoration={brace,amplitude=6pt},thick,softgray]
      ([xshift=10pt]plcp.north east) -- ([xshift=10pt]pmd.south east)
      node[midway,right=9pt,align=left,font=\small,text=coreorange]
      {\textbf{PHY: rebuilt}\\\textbf{every generation}};
  \end{tikzpicture}

  \vspace{0.25em}
  \plainterm{PPDU}{= the PHY-layer transmission unit placed on air:
    preamble + header + payload.}

  \vspace{0.2em}
  \takeaway{Same family behavior above, different ways to turn\\
    bits into radio energy below. That split is Wi-Fi's longevity.}

  % PLCP/PMD is the classic 1997-era split; 802.11-2016+ folds both into a unified
  % PHY clause, and PPDU now expands to "PHY protocol data unit".
  % "association" glossed here at first use; it unblocks 11be's "one association".
  \glossline{MAC = Medium Access Control \quad PLCP = Physical Layer Convergence
    Procedure \quad PMD = Physical Medium Dependent \quad association = joining an
    AP's network \quad (1997 split; later revisions fold PLCP/PMD into one PHY)}
  \source{https://standards.ieee.org/ieee/802.11/10548/}{IEEE Std 802.11 (MAC and PHY specifications)}
\end{frame}

% Teaching note (240 s): problem first. Ethernet detects collisions mid-transmission;
% ask why a radio cannot. Two answers: your own transmitter drowns your receiver, and
% you cannot hear every station anyway (hidden terminals). So 802.11 avoids instead of
% detecting: sense, wait DIFS, random backoff, transmit, and demand a positive ACK
% because silence proves nothing. Optional RTS/CTS handles the hidden-station case.
% The 1997 rates (1-2 Mb/s) look quaint; the access discipline is still in Wi-Fi 7.
\begin{frame}{1997 solved shared access before speed}
  \ask{Ethernet detects collisions mid-send. Why can't a radio?}

  \vspace{0.1em}
  \centering
  % SOURCE: https://standards.ieee.org/wp-content/uploads/interactive/web/wi-fi-timeline/index.html
  % SOURCE: https://standards.ieee.org/ieee/802.11/10548/
  \begin{tikzpicture}[scale=0.88,transform shape]
    \node[netnode,minimum width=1.35cm] (s) at (0,0) {sense};
    \node[netnode,minimum width=1.35cm] (d) at (1.95,0) {DIFS};
    \node[netnode,minimum width=1.5cm] (b) at (3.95,0) {backoff};
    \node[netnode,minimum width=1.35cm] (t) at (5.95,0) {DATA};
    \node[netnode,minimum width=1.35cm,draw=softgreen,fill=green!7] (a) at (7.9,0) {ACK};
    \draw[flow] (s) -- (d); \draw[flow] (d) -- (b);
    \draw[flow] (b) -- (t); \draw[flow] (t) -- (a);
    \node[font=\scriptsize,text=softgray] at (3.95,-0.72)
      {DIFS (Distributed Inter-Frame Space) = a fixed polite pause after the channel goes quiet};
  \end{tikzpicture}

  \vspace{0.15em}
  \begin{columns}[T,onlytextwidth]
    \column{0.42\textwidth}
      \textbf{Original PHY choices}
      \vspace{0.2em}
      \small
      \begin{itemize}\itemsep2pt
        % SOURCE: https://standards.ieee.org/ieee/802.11/10548/ (802.11-1997 base standard)
        \item 1 or 2 Mb/s
        \item FHSS or DSSS at 2.4 GHz
        \item Infrared (never shipped)
      \end{itemize}
    \column{0.54\textwidth}
      \textbf{Because a radio\ldots}
      \vspace{0.2em}
      \small
      \begin{itemize}\itemsep2pt
        \item Own TX deafens own RX
        \item Not every station is audible
        \item So: back off, then demand ACK
      \end{itemize}
  \end{columns}

  \vspace{0.1em}
  \takeaway{CSMA/CA avoids collisions; it cannot detect them\\
    while transmitting. Silence proves nothing --- only an ACK does.}

  \glossline{CSMA/CA = Carrier Sense Multiple Access / Collision Avoidance \quad
    FHSS/DSSS = Frequency-Hopping / Direct-Sequence Spread Spectrum}
  \source{https://standards.ieee.org/ieee/802.11/10548/}{IEEE Std 802.11-1997 baseline}
\end{frame}
